\section{Conclusion and Future Directions}
\label{sec:con}

We have systematically reviewed existing applications, approaches, and evaluation methods for extractive KG summarization. This interdisciplinary research topic spanning artificial intelligence, data mining, and information retrieval provides increasing research opportunities. We highlight the following research directions that are currently underexplored.

\paragraph{Neural Extraction}
While existing approaches mainly exploit the symbolic features of KGs (e.g., ontological schema, graph structure), there is much room for exploring neural methods, given the rapid advances in deep learning models in recent years, particularly graph neural networks.

\paragraph{Supervised Extraction}
Although unsupervised extraction has become mainstream to KG summarization, supervised methods deserve to be taken into consideration. The inadequacy of labeled data for training could be addressed by semi-supervised learning such as self-training and co-training.

\paragraph{Generative Extraction}
Extractive summaries are not necessarily extracted. Motivated by the recent success in generative information extraction and generative information retrieval, it would be an interesting investigation to apply pretrained generative models to KG summary extraction.

\paragraph{Comparative Extraction}
When multiple KGs need to be summarized and compared, e.g., in the results page of a KG search engine, the extraction of their summaries should not be independent but can be performed jointly, attending to their common and distinctive features to facilitate user's selection.

\paragraph{Collaborative Extraction}
Given the computational cost of processing a large KG, we anticipate a scenario where summaries are extracted not from scratch but building on stored summaries previously extracted for similar tasks, which can be inexpensively assembled and tailored to new needs.